\hypertarget{interview-guides}{%
\chapter{Interview Guides}\label{interview-guides}}

\hypertarget{group-a-interview-guide-ioc-technical-staff-and-management}{%
\section{Group A Interview Guide (IOC Technical Staff and
Management)}\label{group-a-interview-guide-ioc-technical-staff-and-management}}

COUHES Protocol E-7658

Estimated Duration: 45-60 minutes

Role-Specific Probe Tracks

This guide contains standardized core questions asked identically to
every Group A participant. Probes are marked with role indicators to
help the interviewer follow the most productive knowledge path for each
participant:

{[}T{]} = Technical/Operational roles (engineers, production managers,
HSE managers, wellsite supervisors, operations supervisors). These
probes ask about operational practices, equipment, procedures, and what
happened at the asset level.

{[}M{]} = Management/Strategic roles (asset managers, country managers,
business development managers, corporate strategy staff). These probes
ask about corporate decisions, resource allocation, strategic rationale,
and organizational design.

{[}BOTH{]} = Applicable to all participants regardless of role.

The interviewer should use all {[}BOTH{]} probes and the probes matching
the participant's role. If a participant's role spans both categories,
the interviewer uses judgment on which probes are within the
participant's firsthand knowledge.

Before You Begin (5 minutes)

\begin{enumerate}
\def\labelenumi{\arabic{enumi}.}
\item
  Introduce yourself and the study: `I am Ikenna Ekekwe, a graduate
  student at MIT. I am conducting research on the operational outcomes
  of IOC asset divestments in Nigeria for my master's thesis. This
  research is not sponsored by or subject to approval by Chevron or any
  other oil company.'
\item
  Disclose your Chevron affiliation: `In the interest of transparency, I
  should mention that I am currently employed by Chevron Nigeria. The
  views in this thesis are my own and independent of my employer.'
\item
  Confirm consent: `Before we begin, I want to confirm that you received
  the study information, that you understand the purpose of this
  interview, and that you agree to participate voluntarily. Do you
  consent to participate?'
\item
  Confirm recording preference: `Do you consent to audio recording of
  this interview, or would you prefer that I take written notes only?'
\item
  Remind participant: `You may decline to answer any question and may
  stop the interview at any time. Your responses will be de-identified.
  You will not be named in the thesis.'
\end{enumerate}

Opening Question (5 minutes)

A1. Please describe your role and responsibilities during the period
when {[}the specific asset/OML{]} was being operated and subsequently
divested.

Probe A1.1 {[}BOTH{]}: What specific operations or functions were you
responsible for?

Probe A1.2 {[}BOTH{]}: Over what time period were you involved?

Probe A1.3 {[}BOTH{]}: How many people reported to you or worked
alongside you in this function?

Primary Questions (35-45 minutes)

Pre-Divestment Capability Architecture

A2. Before the divestment, what operational systems, processes, and
personnel arrangements were in place to run {[}the specific asset{]}?

Probe A2.1 {[}T{]}: Walk me through what happened on a typical day at
the flowstation or production facility.

Probe A2.2 {[}T{]}: When a well integrity issue or equipment failure
arose, what was the response process? Who was involved?

Probe A2.3 {[}T{]}: How were new engineers or operators trained? What
did the training look like in practice?

Probe A2.4 {[}BOTH{]}: Were there operational procedures, knowledge, or
systems that existed only inside the organization, in people's heads or
in informal practices, and were not written down in documents that would
go to a buyer?

Probe A2.5 {[}M{]}: How was the operational capability of this asset
resourced and maintained at the organizational level? What investment,
staffing, and management attention did it require to keep running at the
standard the company maintained?

Pre-Divestment Deterioration

A3. In the years leading up to the divestment, and particularly between
the announcement and the actual handover, did you observe any changes in
how the asset was operated?

Probe A3.1 {[}BOTH{]}: When did you first notice a shift in the
company's commitment to the onshore portfolio? Was this before the
formal divestment announcement?

Probe A3.2 {[}T{]}: Were staff reassigned to other operations
(deepwater, LNG) before the transaction closed? Which functions lost
people first?

Probe A3.3 {[}T{]}: Were maintenance programs, capital investments, or
training activities deferred, reduced, or discontinued during this
period? Were the cuts uniform across all operational functions, or were
some categories, such as training, mentoring, or knowledge transfer
activities, reduced earlier or more severely than production-related
expenditure?

Probe A3.4 {[}T{]}: Were there wells, flowlines, or facilities that were
left in a condition that would require significant intervention by the
incoming operator?

Probe A3.5 {[}M{]}: Was the reduction in investment, staffing, or
maintenance a deliberate corporate decision linked to the planned
divestment, or was it driven by other factors such as cost-cutting or
portfolio rebalancing? Who made these decisions, and what was the
rationale?

Staff Decisions

A4. What happened to the technical and operational staff when the asset
was divested?

Probe A4.1 {[}M{]}: Were staff offered employment by the acquiring
company? If so, on what terms? If not, why not?

Probe A4.2 {[}M{]}: Did the IOC offer to second or loan staff to the
buyer for a transition period? What was the buyer's response?

Probe A4.3 {[}T{]}: Which staff left before the transaction closed
because they saw the divestment coming? Which functions were affected
most?

Probe A4.4 {[}BOTH{]}: Of the staff who had the deepest operational
knowledge of the asset, how many stayed with the new operator and how
many left?

Transfer Process

A5. What specifically was included in the asset transfer when the
transaction closed?

Probe A5.1 {[}T{]}: Were operating procedures, maintenance schedules,
and HSE protocols transferred as documents? Were they complete?

Probe A5.2 {[}BOTH{]}: Was there a formal transition period during which
IOC staff worked alongside the new operator? How long? What did it
cover?

Probe A5.3 {[}BOTH{]}: Were there systems, processes, or institutional
knowledge that you expected to be transferred but were not? What
prevented the transfer?

Probe A5.4 {[}M{]}: What information did the buyer have access to about
how the asset was actually operated, beyond the formal data room
documents? Could the buyer assess the operational capability it was
acquiring before the transaction closed?

Probe A5.5 {[}M{]}: Who decided what would be included in the data room
and the transfer package? Were there categories of operational knowledge
or organizational capability that were deliberately excluded? If so,
what was the rationale?

Probe A5.6 {[}M{]}: Were there pressures within the company to ensure a
responsible or safe transition, such as reputational concerns,
regulatory expectations, or liability considerations? How were these
balanced against the commercial priorities of the divestment?

Contractor and Vendor Ecosystem

A6. How did the relationships with service companies, drilling
contractors, and maintenance vendors work, and what happened to those
relationships after the divestment?

Probe A6.1 {[}T{]}: Which service companies were critical to the
day-to-day operation of the asset? What did they provide?

Probe A6.2 {[}BOTH{]}: Were these contractor relationships tied to the
IOC's corporate agreements, or were they asset-specific?

Probe A6.3 {[}M{]}: Would the same contractors extend the same terms,
priority, and crew quality to the new operator? What would you expect to
change?

Probe A6.4 {[}T{]}: Were there specialized capabilities (well control,
subsea intervention, corrosion management) that depended entirely on
contractor relationships?

Community and Stakeholder Architecture

A7. How did the relationships with host communities, security
arrangements, and stakeholder agreements work, and what happened to them
during the divestment?

Probe A7.1 {[}M{]}: Were there community agreements (Global Memoranda of
Understanding, MOUs) that were tied to the IOC rather than the asset?
Did they transfer?

Probe A7.2 {[}T{]}: Were there community liaison officers or other staff
whose relationships with community leaders were critical to operational
access?

Probe A7.3 {[}M{]}: How were security arrangements (military,
community-based, private) organized, and did they transfer with the
asset?

Probe A7.4 {[}BOTH{]}: Did you anticipate any community-related
operational disruptions for the incoming operator?

Coordination from the Seller's Perspective

A8. How was the transition coordinated among the seller, the buyer, and
the regulator?

Probe A8.1 {[}M{]}: Was there a formal transition plan? Who designed it?
What did it cover?

Probe A8.2 {[}BOTH{]}: At what point did coordination break down, if it
did: before the transaction was approved, during the handover period, or
after the new operator took control?

Probe A8.3 {[}BOTH{]}: Did the regulator play any role in overseeing or
verifying the transition? What specifically did the regulator do or not
do?

Probe A8.4 {[}BOTH{]}: Looking back, what would have needed to be
coordinated differently for the transition to produce a better
operational outcome?

Probe A8.5 {[}M{]}: To your knowledge, what did the regulator assess
about the acquiring company before approving the transaction? Were there
aspects of operational readiness that you believe should have been
assessed but were not?

Secondary Question: JV Governance

A9. How did the joint venture governance with NNPC work, and how would
you expect it to change with a new operator?

Probe A9.1 {[}M{]}: How were cash calls, budget approvals, and work
programs managed within the JV?

Probe A9.2 {[}BOTH{]}: Was NNPC's engagement with the IOC operator
different from what you would expect with a domestic operator?

Framework Elicitation

A10. If you could redesign the divestment process to prevent the
operational failures that have occurred, what would you change?

Probe A10.1 {[}BOTH{]}: What should the seller be required to do before
or during the transfer?

Probe A10.2 {[}BOTH{]}: What should the buyer be required to demonstrate
before being approved to acquire the asset?

Probe A10.3 {[}BOTH{]}: What should the regulator be required to verify
before, during, and after the transaction?

Probe A10.4 {[}BOTH{]}: Are there examples of transfers that worked
well? What made them different?

Closing (5 minutes)

\begin{enumerate}
\def\labelenumi{\arabic{enumi}.}
\item
  `Is there anything important about the divestment process or its
  consequences that my questions did not cover?'
\item
  `Is there anyone else you would recommend I speak with who has direct
  knowledge of these events?'
\item
  `Do you have any questions about the study?'
\item
  Thank the participant. Confirm their preferred contact method if
  follow-up questions arise.
\item
  Immediately after the call, complete expanded notes while the
  conversation is fresh. Note any statements the participant emphasized
  or that you judge to be analytically significant. Record the
  participant's role type (technical/operational or
  management/strategic) for analytical coding.
\end{enumerate}

\hypertarget{group-b-interview-guide-domestic-oil-company-operators-and-managers}{%
\section{Group B Interview Guide (Domestic Oil Company Operators and
Managers)}\label{group-b-interview-guide-domestic-oil-company-operators-and-managers}}

Group B: DOC Operators and Managers

COUHES Protocol E-7658

Estimated Duration: 45-60 minutes

Role-Specific Probe Tracks

This guide contains standardized core questions asked identically to
every Group B participant. Probes are marked with role indicators:

{[}T{]} = Technical/Operational roles (production engineers, field
operators, HSE staff, maintenance engineers). These probes ask about
what happened operationally, what gaps were encountered, and what the
participant experienced at the asset level.

{[}M{]} = Management/Strategic roles (managing directors, business
development staff, acquisition team leaders, CFOs, operations
directors). These probes ask about strategic decisions, financial
positions, organizational design, and the rationale behind key choices.

{[}BOTH{]} = Applicable to all participants regardless of role.

The interviewer should use all {[}BOTH{]} probes and the probes matching
the participant's role.

Before You Begin (5 minutes)

\begin{enumerate}
\def\labelenumi{\arabic{enumi}.}
\item
  Introduce yourself and the study: `I am Ikenna Ekekwe, a graduate
  student at MIT. I am conducting research on the operational outcomes
  of IOC asset divestments in Nigeria for my master's thesis. This
  research is not sponsored by or subject to approval by Chevron or any
  other oil company.'
\item
  Disclose your Chevron affiliation: `In the interest of transparency, I
  should mention that I am currently employed by Chevron Nigeria. The
  views in this thesis are my own and independent of my employer.'
\item
  Confirm consent: `Before we begin, I want to confirm that you received
  the study information, that you understand the purpose of this
  interview, and that you agree to participate voluntarily. Do you
  consent to participate?'
\item
  Confirm recording preference: `Do you consent to audio recording of
  this interview, or would you prefer that I take written notes only?'
\item
  Remind participant: `You may decline to answer any question and may
  stop the interview at any time. Your responses will be de-identified.
  You will not be named in the thesis.'
\end{enumerate}

Opening Question (5 minutes)

B1. Please describe your role and responsibilities during and after the
acquisition of {[}the specific asset/OML{]}.

Probe B1.1 {[}BOTH{]}: What specific operations or functions were you
responsible for?

Probe B1.2 {[}BOTH{]}: Were you involved in the acquisition process
itself, the transition, or only in post-acquisition operations?

Probe B1.3 {[}BOTH{]}: How long after the transaction closed did you
remain involved with the asset?

Primary Questions (35-45 minutes)

Staffing Decisions and Organizational Preparation

B2. What technical and operational staff did your organization bring to
the asset, and how were staffing decisions made?

Probe B2.1 {[}M{]}: Did your organization attempt to retain the IOC's
technical staff? Were you offered the option? If you declined, why?

Probe B2.2 {[}BOTH{]}: Where did your operational staff come from? Were
they experienced in upstream operations at this scale and complexity?

Probe B2.3 {[}T{]}: Were there specific technical functions (well
integrity, HSE, pipeline integrity, emergency response) where your team
lacked experience?

Probe B2.4 {[}M{]}: Before the acquisition, what experience did your
organization have in operating upstream assets at the scale and
complexity of what you were acquiring?

Probe B2.5 {[}M{]}: What was your organization's strategy for building
the operational capability to run this asset? Was there a deliberate
plan to develop capability internally, acquire it through hiring, or
rely on the transition from the IOC? How was this decision made?

Probe B2.6 {[}M{]}: When your organization took over, what operational
functions were established in your organizational structure? Were there
functions that the IOC had maintained, such as well control, emergency
response, or corrosion management, that your organization did not
include in its initial structure?

Expectation vs.~Reality

B3. What did you expect to find when you took over the asset, and how
did the reality compare?

Probe B3.1 {[}M{]}: Based on due diligence and the data room, what was
your understanding of the asset's condition and operational
requirements?

Probe B3.2 {[}T{]}: When you took over, what surprised you? What was
different from what the pre-acquisition information led you to expect?

Probe B3.3 {[}T{]}: Were there specific problems (well conditions,
pipeline integrity, environmental liabilities, staffing gaps) that were
worse than disclosed?

Probe B3.4 {[}M{]}: Looking back, was there information you needed
before the acquisition that you did not have and could not have obtained
from the data room?

Post-Divestment Operational Reality

B4. After the transaction closed, what specific operational changes
occurred in the first 6 to 12 months?

Probe B4.1 {[}T{]}: What was the first major operational challenge your
organization encountered?

Probe B4.2 {[}T{]}: Were there specific functions or capabilities that
your organization found it could not perform without outside assistance?

Probe B4.3 {[}M{]}: Did your organization bring in external contractors
or consultants? For what functions? For how long? Who made this
decision? Were external specialists needed to identify what had gone
wrong, or did your organization know what the problem was and bring in
specialists to help execute the response?

Probe B4.4 {[}T{]}: Were there differences between your organization's
ability to handle day-to-day production operations and its ability to
respond to emergencies, equipment failures, or well control situations?
Can you give me specific examples of each?

Probe B4.5 {[}BOTH{]}: Were there any incidents, near-misses, or
operational disruptions during this period? What happened, and how were
they handled?

Probe B4.6 {[}M{]}: Before the first major operational incident or
crisis, had your organization taken any steps to build capability in
that specific area, such as hiring specialists, engaging contractors, or
conducting training? If not, was the need recognized before the event
occurred?

Probe B4.7 {[}BOTH{]}: Over the first two to three years of your
operations, did your organization's capability in any specific area
improve? If so, what caused the improvement: additional funding, hiring
of experienced personnel, organizational changes, or something else? Can
you identify specific changes that made a difference?

Probe B4.8 {[}M{]}: Looking at the period from the first year of
operations to now, how has your organization's operational capability
changed? Has it improved, stayed the same, or declined? If it improved,
what drove the improvement? If it did not improve despite sustained
production, why not?

Financial Resources

B5. What role, if any, did financial resources play in your
organization's ability to operate the asset after the transaction?

Probe B5.1 {[}M{]}: Did your organization have the capital to maintain
the asset's infrastructure, including wells, flowstations, pipelines,
and export facilities?

Probe B5.2 {[}M{]}: Were there specific investments that were needed but
could not be funded? What were they?

Probe B5.3 {[}T{]}: Were there cases where the money was available but
the organization could not execute the work? What prevented execution?

Probe B5.4 {[}M{]}: Was the JV funding (NNPC cash calls) maintained
throughout the period? Were there cash call disputes?

PIA Fiscal Effects

B6. Did the fiscal terms of operating the asset change after the
Petroleum Industry Act of 2021, and if so, how did that affect
operations?

Probe B6.1 {[}M{]}: Were operational decisions affected by the PIA's
fiscal restructuring?

Probe B6.2 {[}BOTH{]}: For assets divested before 2021, were the
challenges you observed related to the fiscal environment or to
something else?

Probe B6.3 {[}M{]}: Did the fiscal terms make a difference in your
assessment between operators who succeeded and those who did not?

Coordination from the Buyer's Perspective

B7. How were transition activities coordinated among the seller, the
buyer, and the regulator?

Probe B7.1 {[}M{]}: Was there a formal transition plan? Who designed it?
What did it cover?

Probe B7.2 {[}BOTH{]}: What transition support did you receive from the
seller? What did you expect but did not receive?

Probe B7.3 {[}BOTH{]}: At what point did coordination break down, if it
did: before the transaction was approved, during the handover period, or
after your organization took full control?

Probe B7.4 {[}BOTH{]}: Did the regulator play any role in overseeing or
verifying the transition? What specifically did the regulator do or not
do?

Probe B7.5 {[}BOTH{]}: Looking back, what would have needed to be
coordinated differently for the transition to produce a better outcome?

Secondary Questions

Security and Oil Theft

B8. Did the security situation or the level of oil theft change after
the IOC departed, and if so, how did that affect your operations?

Probe B8.1 {[}BOTH{]}: Were there changes in community relations,
pipeline vandalism, or militant activity after the operator changed?

Probe B8.2 {[}M{]}: Did the security arrangements that the IOC had in
place transfer with the asset, or did your organization have to rebuild
them?

Probe B8.3 {[}BOTH{]}: To what extent did security-related losses
(theft, vandalism, shut-ins) account for production decline versus
operational capability issues?

Post-Close Regulatory Engagement

B9. After the transaction closed, what role did the regulator play in
overseeing your operations?

Probe B9.1 {[}M{]}: Were there regulatory inspections, compliance
audits, or enforcement actions during the transition period?

Probe B9.2 {[}BOTH{]}: Did the regulator provide any guidance, support,
or oversight during the first 12 months of your operations?

Probe B9.3 {[}T{]}: Was the regulatory environment different for your
organization as a domestic operator compared to what the IOC had
experienced?

JV/NNPC Relationship

B10. How did the joint venture governance with NNPC work after you took
over as operator?

Probe B10.1 {[}M{]}: Were cash calls, budget approvals, and work
programs managed smoothly, or were there disputes?

Probe B10.2 {[}BOTH{]}: Was NNPC's level of engagement or cooperation
different from what the IOC had experienced?

Framework Elicitation

B11. If you could redesign the divestment process to prevent the
operational failures that have occurred, what would you change?

Probe B11.1 {[}BOTH{]}: What should the seller be required to do before
or during the transfer?

Probe B11.2 {[}BOTH{]}: What should the buyer be required to demonstrate
before being approved to acquire the asset?

Probe B11.3 {[}BOTH{]}: What should the regulator be required to verify
before, during, and after the transaction?

Probe B11.4 {[}BOTH{]}: Are there examples of transfers that worked
well? What made them different?

Closing (5 minutes)

\begin{enumerate}
\def\labelenumi{\arabic{enumi}.}
\item
  `Is there anything important about the divestment process or its
  consequences that my questions did not cover?'
\item
  `Is there anyone else you would recommend I speak with who has direct
  knowledge of these events?'
\item
  `Do you have any questions about the study?'
\item
  Thank the participant. Confirm their preferred contact method if
  follow-up questions arise.
\item
  Immediately after the call, complete expanded notes while the
  conversation is fresh. Note any statements the participant emphasized
  or that you judge to be analytically significant. Record the
  participant's role type (technical/operational or
  management/strategic) for analytical coding.
\end{enumerate}

\hypertarget{group-c-interview-guide-regulatory-officials-and-industry-observers}{%
\section{Group C Interview Guide (Regulatory Officials and Industry
Observers)}\label{group-c-interview-guide-regulatory-officials-and-industry-observers}}

COUHES Protocol E-7658. Estimated Duration: 45-60 minutes.

Before You Begin (5 minutes): (1) Introduce yourself and the study. (2)
Disclose Chevron affiliation with explicit statement that the thesis is
not sponsored by or subject to approval by Chevron or any other oil
company. (3) Confirm consent. (4) Confirm recording preference. (5)
Remind participant of right to decline any question and stop the
interview at any time.

C1. Please describe your role and responsibilities in relation to
upstream petroleum operations and IOC asset divestments in Nigeria.

Probe C1.1: What specific regulatory, oversight, or advisory functions
were you responsible for?

Probe C1.2: Over what time period were you involved in
divestment-related activities?

Probe C1.3: Approximately how many divestment transactions were you
involved in or had oversight of?

C2. When an IOC proposed to divest an upstream asset to a domestic
company, what did the regulatory approval process involve?

Probe C2.1: What specific criteria were used to evaluate whether the
acquiring company was qualified to take over the asset?

Probe C2.2: Was the acquiring company's operational or technical
capability assessed as part of the approval process? If so, how? If not,
why not?

Probe C2.3: Was there a legal basis in the Petroleum Act, the PIA, or
other legislation that would have allowed the regulator to require a
demonstration of operational capability as a condition of approval?

Probe C2.4: Was there anything the regulator wanted to assess but lacked
the legal authority or the institutional capacity to require? What was
it, and what prevented it?

Probe C2.5: Was the buyer's financial capacity to maintain and operate
the asset assessed as part of the approval process? In cases where the
buyer subsequently experienced financial difficulties, was there any
indication of this at the time of approval?

Probe C2.6: When assessing a buyer's qualification to operate the asset,
was there a distinction between the buyer's formal qualifications, such
as engineering credentials, corporate structure, and financial capacity,
and the practical operational knowledge required to run the specific
asset? Was the regulator in a position to assess whether the buyer could
handle the operational realities of the specific asset, as opposed to
upstream operations in general?

Probe C2.7: Was any proposed divestment transaction ever rejected
because the acquiring company was judged to lack the operational
capability to run the asset? If not, was that because all buyers were
judged capable, or because operational capability was not a criterion on
which a transaction could be rejected?

C3. Divestment approval involves ministerial consent under Section 29 of
the Petroleum Act. How did the ministerial consent process interact with
the regulatory assessment?

Probe C3.1: Were there cases where the technical or regulatory
assessment raised concerns about a proposed buyer, but the transaction
was approved at the ministerial level?

Probe C3.2: To what extent did considerations other than operational
capability, such as indigenous ownership targets, political
relationships, or commercial terms, influence divestment approvals?

Probe C3.3: In your experience, was there a tension between the policy
objective of promoting indigenous ownership and the objective of
ensuring operational continuity?

C4. After a divestment was approved and the new operator took over, what
regulatory oversight was in place?

Probe C4.1: Were there scheduled inspections, compliance audits, or
performance reviews of the new operator during the first 12 to 24
months?

Probe C4.2: When post-divestment operational problems emerged
(production decline, safety incidents, environmental pollution), what
enforcement tools were available to the regulator? Were they used?

Probe C4.3: Was there a formal mechanism for the regulator to intervene
if the new operator's performance fell below acceptable standards?

Probe C4.4: Was the level of regulatory oversight different for newly
divested assets compared to established operations?

C5. How has the Petroleum Industry Act of 2021 changed the regulatory
framework for upstream asset divestments?

Probe C5.1: Did the PIA introduce new requirements for assessing buyer
capability, or does the regulatory gap that existed under the Petroleum
Act persist?

Probe C5.2: How has the creation of the NUPRC changed the institutional
capacity for overseeing divestments compared to the former DPR?

Probe C5.3: In your assessment, did the PIA's fiscal restructuring
contribute to the operational challenges observed among domestic
operators, or were the challenges primarily driven by other factors?

C6. How was the transition process coordinated among the seller, the
buyer, and the regulator?

Probe C6.1: Did the regulator have a defined role in overseeing or
facilitating the transition from the IOC to the new operator?

Probe C6.2: Were there cases where the regulator attempted to coordinate
transition activities? What happened?

Probe C6.3: At what point in the process did coordination typically
break down, if it did: before approval, during handover, or after close?

Probe C6.4: Was there a regulatory mechanism for ensuring that the
operational capability to run the asset would survive the transaction?

C7. Across the divestment transactions you were involved in or observed,
did you notice systematic differences in how acquiring companies
performed after taking over?

Probe C7.1: Which acquisitions went well and which did not? In your
assessment, what distinguished the successful ones from the unsuccessful
ones?

Probe C7.2: Were there common patterns in the types of operational
problems that emerged after divestments?

Probe C7.3: Were some types of buyers (those with prior operating
experience, international partnerships, IOC staff retention)
systematically more successful than others?

Probe C7.4: Does the pattern you observed suggest that the problem is
primarily about individual operator quality, or about something
structural in how divestments are designed and approved?

Probe C7.5: Did you observe parallels between the outcomes of the
marginal field program (2003 onwards) and the subsequent IOC divestment
outcomes? Were lessons from the marginal field experience considered
when the IOC divestment framework was applied?

C8. The NUPRC announced a seven-pillar divestment assessment framework
in 2024. Can you describe what prompted the development of this
framework and what it is designed to address?

Probe C8.1: What evidence or experience led the NUPRC to conclude that a
new framework was needed?

Probe C8.2: What specific gaps in the previous regulatory approach is
the framework designed to close?

Probe C8.3: Does the framework include requirements for the seller to
disclose operational capability, for the buyer to demonstrate readiness,
or for the regulator to verify capability transfer?

Probe C8.4: How does the new framework compare to what existed before,
and what do you expect its impact to be?

C9. If you could redesign the divestment process to prevent the
operational failures that have occurred, what would you change?

Probe C9.1: What should the seller be required to do before or during
the transfer?

Probe C9.2: What should the buyer be required to demonstrate before
being approved to acquire the asset?

Probe C9.3: What should the regulator be required to verify before,
during, and after the transaction?

Probe C9.4: Should capability transfer be a mandatory condition of
divestment approval, or should it remain at the discretion of the
transacting parties?

Closing (5 minutes): (1) Is there anything important about the
divestment process or its consequences that my questions did not cover?
(2) Is there anyone else you would recommend I speak with who has direct
knowledge of these events? (3) Do you have any questions about the
study? (4) Thank the participant. Confirm preferred contact method for
follow-up. (5) Immediately after the call, complete expanded notes while
the conversation is fresh.
